\documentclass[11pt]{article}
\usepackage[a4paper,margin=27mm]{geometry}
\usepackage{amsmath,amssymb,mathtools}
\usepackage{fontspec,unicode-math}
\setmainfont{Latin Modern Roman}
\setmathfont{Latin Modern Math}
\usepackage{xurl}
\usepackage[hidelinks]{hyperref}
\newcommand{\sref}[2]{\hyperlink{source:#1:#2}{[S#2]}}
\newcommand{\cataloguescope}[1]{\paragraph{Mathematical question.}#1}
\hypersetup{pdftitle={UnsolvedMath Problem 9400151}}
\begin{document}
% UnsolvedMath ID 9400151; source catalogue code AMR-093-0151
\setcounter{section}{9400150}
\section{Are there infinitely many Kynea primes}\label{9400151}
{\small\sffamily UnsolvedMath ID 9400151 · Number Theory}
\subsection{Definitions and mathematical statement}

\paragraph{Shared notation.}$\mathbb N=\{1,2,\ldots\}$, and $\mathbb Z,\mathbb Q,\mathbb R,\mathbb C$ denote the integers, rationals, reals and complex numbers. Any other conventions are specified locally.


For integers $n\ge0$, define $K_n=(2^n+1)^2-2=4^n+2^{n+1}-1$. A Kynea prime is a prime value $K_n$.
\paragraph{Statement recorded in the supplied source.}
Are there infinitely many $n\ge0$ for which $K_n$ is prime? \sref{9400151}{2}
\subsection{Short English statement}
Are infinitely many Kynea numbers prime?
\subsection{Sources}
\begin{description}
\item[\hypertarget{source:9400151:1}{S1}] ULAM.AI and the UnsolvedMath Contributors, UnsolvedMath: A Curated Collection of Open Mathematics Problems, record 9400151 (AMR-093-0151). CC BY 4.0. \url{https://huggingface.co/datasets/ulamai/UnsolvedMath}.
\item[\hypertarget{source:9400151:2}{S2}] Wikipedia contributors, \emph{List of unsolved problems in mathematics}, revision 1366636767 (2026) Number theory, Are there infinitely many Kynea primes. \url{https://en.wikipedia.org/w/index.php?title=List_of_unsolved_problems_in_mathematics&oldid=1366636767#Number_theory}.
\item[\hypertarget{source:9400151:3}{S3}] OEIS Foundation and contributors, \emph{The On-Line Encyclopedia of Integer Sequences}, A091513: Numbers k such that (2\textasciicircum{}k + 1)\textasciicircum{}2 - 2 = 4\textasciicircum{}k + 2\textasciicircum{}(k+1) - 1 is prime., official data snapshot \texttt{693b5499c9d2} (2026) Name, definition and comments. \url{https://raw.githubusercontent.com/oeis/oeisdata/693b5499c9d248f570dd2d979e8bd1d7eecf689b/seq/A091/A091513.seq}.
\end{description}
\label{end:9400151}
\end{document}
